% Teil II: durchgerechneter Weg, Teile a--c.

\subsection*{Gegeben und gesucht}
Masse $m=\qty{100}{g}=0{,}1\ \unit{kg}$, Federkonstante
$k=\qty{40}{N/m}$, Reibung vernachlaessigbar. Gesucht: (a) Ruhelage $z_0$,
(b) Bewegungsgleichung, allgemeine Loesung, $\omega$, $T$,
(c) Anfangsbedingungen und $z(t)$ fuer das Loslassen aus $30\,\mathrm{cm}$.

\subsection*{Ansatz}
Ruhelage aus dem Kräftegleichgewicht; mit $u=z-z_0$ erhaelt man den
harmonischen Oszillator; die Loesungskonstanten aus den
Anfangsbedingungen bestimmen.

\subsection*{Rechnung}
\textbf{a) Ruhelage.} Die Federkraft $k z_0$ (hier: $z$ nach unten von der
unbelasteten Feder) haelt der Gewichtskraft: $k z_0=m g$, also
\begin{equation}
  z_0=\frac{m g}{k}
  =\frac{\qty{0.1}{kg}\cdot\qty{9.81}{m/s^2}}{\qty{40}{N/m}}
  =0{,}0245\ \mathrm m=\qty{2.45}{cm}\;.
\end{equation}

\textbf{b) Bewegungsgleichung und allgemeine Loesung.} Fuer die Auslenkung
$z$ relativ zur Ruhelage $z_0$ ist die Federkraft $F=-k(z-z_0)$; Newton:
\begin{equation}
  m\ddot z=-k\,(z-z_0)
  \;\Rightarrow\;
  m\ddot u+k u=0,\quad u=z-z_0,\qquad
  \ddot u+\omega^2u=0
\end{equation}
mit
\begin{equation}
  \omega=\sqrt{\frac{k}{m}}=\sqrt{\frac{40}{0.1}}\approx\qty{20}{s^{-1}},
  \qquad
  T=\frac{2\pi}{\omega}\approx\qty{0.314}{s}\;.
\end{equation}
Allgemeine Loesung (mit frei waehlbaren $A,B$):
\begin{equation}
  u(t)=A\cos(\omega t)+B\sin(\omega t),\qquad
  \dot u(t)=-\omega A\sin(\omega t)+\omega B\cos(\omega t).
\end{equation}

\textbf{c) Anfangsbedingungen und Bewegung.} Als Ursprung ist die Ruhelage
($z=0$) definiert; $z$ z\"ahlt nach oben positiv. Die Kugel wird \emph{30 cm
nach unten} gedrueckt und bei $t=0$ aus der Ruhe losgelassen:
\begin{equation}
  z(0)=-0{,}30\,\mathrm m,\qquad \dot z(0)=0\quad(\text{losgelassen}).
\end{equation}
Damit ist $u(0)=-0{,}30\,\mathrm m$, $\dot u(0)=0$; aus $u(0)=A$ und
$\dot u(0)=\omega B=0$ folgt $A=-0{,}30\,\mathrm m$, $B=0$:
\begin{equation}
  \boxed{\;z(t)=-0{,}30\,\mathrm m\cdot\cos(\qty{20}{s^{-1}}\,t)\;},\qquad
  \dot z(t)=+\qty{6}{m/s}\,\sin(\qty{20}{s^{-1}}\,t)\;.
\end{equation}
Die Kugel schwingt um die Ruhelage $z=0$ mit Amplitude $\qty{0.30}{m}$,
Periode $T\approx0{,}31\,\mathrm s$; bei $t=0$ ist sie 30\,cm unter der
Ruhelage und startet aus der Ruhe (Skizze: Kosinus nach unten aufgesetzt).

\begin{ergebnis}
(a) $z_0=\qty{2.45}{cm}$; (b) $\omega\approx\qty{20}{s^{-1}}$,
$T\approx\qty{0.314}{s}$, $z(t)=A\cos(\omega t)+B\sin(\omega t)$;
(c) $z(0)=-0{,}30\,\mathrm m$, $\dot z(0)=0$, damit
$z(t)=-0{,}30\,\mathrm m\cdot\cos(\qty{20}{s^{-1}}t)$.
\end{ergebnis}

\begin{kontrolle}
\textbf{Einheiten:} $\omega=\sqrt{k/m}$: $\sqrt{\mathrm{N/(m\,kg)}}=
\sqrt{(\mathrm{kg\,m/s^2})/(\mathrm{m\,kg})}=\mathrm{s^{-1}}$; $T$ in s;
Argumente $\omega t$ dimensionslos. \checkmark

\textbf{Querprobe:} Bei $t=0$ ist $z=-0{,}30\,\mathrm m$ (maximal nach
unten), $\dot z=0$ (Umkehrpunkt), wie erwartet. $\omega=20\,\mathrm{s^{-1}}$
liefert $T\approx0{,}31\,\mathrm s$, vertraeglich mit typischen Federpendeln.
\checkmark

\textbf{DGL:} Einsetzen $u=A\cos(\omega t)$ in $\ddot u+\omega^2u=0$
erfuellt. \checkmark
\end{kontrolle}

\begin{merksatz}
Fuer eine vertikale Feder rechnet man den Schwingungsursprung in die
Ruhelage $z_0=mg/k$; relative Auslenkungen schwingen mit
$\omega=\sqrt{k/m}$ und $T=2\pi\sqrt{m/k}$. Die Amplitude ist die
Auslenkung zum Loslassen, wenn die Kugel dabei in Ruhe ist.
\end{merksatz}
